\chapter{Modelo de datos}
\label{cap:AnexoD}

Este anexo documenta el esquema de base de datos del sistema. PostgreSQL 15 aloja dos bases de datos con propósitos distintos: \texttt{cie10\_eventstore}, gestionada por la biblioteca Commanded y que actúa como registro de eventos inmutable, y \texttt{cie10\_app}, que contiene las proyecciones de lectura y los datos de referencia del catálogo CIE-10-ES.

\section{Diagrama entidad-relación}

La Figura~\ref{fig:erd-bd} muestra el diagrama entidad-relación de la base de datos \texttt{cie10\_app}, con las tablas de proyección y la tabla de referencia del catálogo CIE-10-ES.

\begin{figure}[ht]
\centering
\resizebox{\textwidth}{!}{% Diagrama entidad-relación — Esquema de base de datos
% Uso: % Diagrama entidad-relación — Esquema de base de datos
% Uso: % Diagrama entidad-relación — Esquema de base de datos
% Uso: \input{figs/erd-bd} dentro de figure+\centering+\resizebox{\textwidth}{!}{...}
%
% Tablas: users, conversations, messages, analysis_cards,
%         predicted_codes, code_suggestions, cie10_codes

\begin{tikzpicture}[
  every node/.style={font=\scriptsize\sffamily},
  %--- Estilo de tabla ERD (cabecera + campos) ---
  tbl/.style={
    rectangle split, rectangle split parts=2,
    rectangle split part fill={aquaESI!70, white},
    draw=aquaESI!80, thick, rounded corners=3pt,
    align=left, text width=3.8cm
  },
  tblsm/.style={
    rectangle split, rectangle split parts=2,
    rectangle split part fill={turquesa!60, white},
    draw=turquesa!70, thick, rounded corners=3pt,
    align=left, text width=3.4cm
  },
  tblref/.style={
    rectangle split, rectangle split parts=2,
    rectangle split part fill={gray!50, white},
    draw=gray!60, thick, rounded corners=3pt,
    align=left, text width=3.4cm
  },
  %--- Flechas de FK (1:N) ---
  fk/.style={->, >=Stealth, thick, gray!60},
  node distance=1.0cm and 1.6cm
]

%--- NODOS ---

% users
\node[tbl] (users) {
  \textbf{users}
  \nodepart{two}
  \underline{id} : uuid PK\\
  first\_name, last\_name\\
  username, email\\
  password\_hash\\
  locale, confirmed\_at\\
  inserted\_at, updated\_at
};

% conversations
\node[tbl, right=2.0cm of users] (conv) {
  \textbf{conversations}
  \nodepart{two}
  \underline{id} : uuid PK\\
  conversation\_id : string\\
  user\_id : uuid FK\\
  started\_at : utc\_datetime\\
  status : string\\
  deleted\_at : utc\_datetime\\
  inserted\_at, updated\_at
};

% messages
\node[tblsm, above right=0.8cm and 1.6cm of conv] (msg) {
  \textbf{messages}
  \nodepart{two}
  \underline{id} : uuid PK\\
  message\_id : string\\
  content : text\\
  user\_id, message\_type\\
  timestamp : utc\_datetime\\
  conversation\_id : uuid FK\\
  inserted\_at, updated\_at
};

% analysis_cards
\node[tblsm, right=1.6cm of conv] (cards) {
  \textbf{analysis\_cards}
  \nodepart{two}
  \underline{id} : uuid PK\\
  card\_id : string\\
  card\_type : string\\
  content : text\\
  position : integer\\
  message\_id : string\\
  conversation\_id : uuid FK\\
  inserted\_at, updated\_at
};

% predicted_codes
\node[tblsm, below right=0.8cm and 1.6cm of conv] (pred) {
  \textbf{predicted\_codes}
  \nodepart{two}
  \underline{id} : uuid PK\\
  code\_id : string\\
  cie10\_code : string\\
  confidence\_score : float\\
  status : string\\
  validated\_by, rejected\_by\\
  conversation\_id : uuid FK\\
  inserted\_at, updated\_at
};

% code_suggestions
\node[tblsm, below=1.0cm of conv] (sugg) {
  \textbf{code\_suggestions}
  \nodepart{two}
  \underline{id} : uuid PK\\
  suggestion\_id : string\\
  selected\_text : text\\
  suggested\_code : string\\
  suggested\_by : string\\
  conversation\_id : uuid FK\\
  inserted\_at, updated\_at
};

% cie10_codes
\node[tblref, right=1.6cm of pred] (cie10) {
  \textbf{cie10\_codes}
  \nodepart{two}
  \underline{id} : uuid PK\\
  code : string UNIQUE\\
  description : text\\
  type : string\\
  metadata : jsonb\\
  inserted\_at, updated\_at
};

%--- RELACIONES FK ---

% users 1 → N conversations
\draw[fk] (conv.west) -- node[above, font=\tiny\sffamily]{N} ++(-0.4,0)
  node[above, font=\tiny\sffamily, at end] {} -- (users.east)
  node[above, font=\tiny\sffamily, pos=0.9]{1};

% conversations 1 → N messages
\draw[fk] (msg.south) -- node[right, font=\tiny\sffamily]{N} ++(0,-0.3)
  -| (conv.north) node[left, font=\tiny\sffamily, pos=0.9]{1};

% conversations 1 → N analysis_cards
\draw[fk] (cards.west) -- node[above, font=\tiny\sffamily]{N} ++(-0.3,0)
  -- (conv.east) node[above, font=\tiny\sffamily, pos=0.1]{1};

% conversations 1 → N predicted_codes
\draw[fk] (pred.north) -- node[right, font=\tiny\sffamily]{N} ++(0,0.3)
  -| (conv.south) node[left, font=\tiny\sffamily, pos=0.9]{1};

% conversations 1 → N code_suggestions
\draw[fk] (sugg.north) -- node[right, font=\tiny\sffamily]{N} ++(0,0.2)
  -- (conv.south) node[right, font=\tiny\sffamily, pos=0.05]{1};

% predicted_codes N → 1 cie10_codes (referencia por código)
\draw[fk] (pred.east) -- node[above, font=\tiny\sffamily]{N} ++(0.3,0)
  -- (cie10.west) node[above, font=\tiny\sffamily, pos=0.9]{1};

\end{tikzpicture}
 dentro de figure+\centering+\resizebox{\textwidth}{!}{...}
%
% Tablas: users, conversations, messages, analysis_cards,
%         predicted_codes, code_suggestions, cie10_codes

\begin{tikzpicture}[
  every node/.style={font=\scriptsize\sffamily},
  %--- Estilo de tabla ERD (cabecera + campos) ---
  tbl/.style={
    rectangle split, rectangle split parts=2,
    rectangle split part fill={aquaESI!70, white},
    draw=aquaESI!80, thick, rounded corners=3pt,
    align=left, text width=3.8cm
  },
  tblsm/.style={
    rectangle split, rectangle split parts=2,
    rectangle split part fill={turquesa!60, white},
    draw=turquesa!70, thick, rounded corners=3pt,
    align=left, text width=3.4cm
  },
  tblref/.style={
    rectangle split, rectangle split parts=2,
    rectangle split part fill={gray!50, white},
    draw=gray!60, thick, rounded corners=3pt,
    align=left, text width=3.4cm
  },
  %--- Flechas de FK (1:N) ---
  fk/.style={->, >=Stealth, thick, gray!60},
  node distance=1.0cm and 1.6cm
]

%--- NODOS ---

% users
\node[tbl] (users) {
  \textbf{users}
  \nodepart{two}
  \underline{id} : uuid PK\\
  first\_name, last\_name\\
  username, email\\
  password\_hash\\
  locale, confirmed\_at\\
  inserted\_at, updated\_at
};

% conversations
\node[tbl, right=2.0cm of users] (conv) {
  \textbf{conversations}
  \nodepart{two}
  \underline{id} : uuid PK\\
  conversation\_id : string\\
  user\_id : uuid FK\\
  started\_at : utc\_datetime\\
  status : string\\
  deleted\_at : utc\_datetime\\
  inserted\_at, updated\_at
};

% messages
\node[tblsm, above right=0.8cm and 1.6cm of conv] (msg) {
  \textbf{messages}
  \nodepart{two}
  \underline{id} : uuid PK\\
  message\_id : string\\
  content : text\\
  user\_id, message\_type\\
  timestamp : utc\_datetime\\
  conversation\_id : uuid FK\\
  inserted\_at, updated\_at
};

% analysis_cards
\node[tblsm, right=1.6cm of conv] (cards) {
  \textbf{analysis\_cards}
  \nodepart{two}
  \underline{id} : uuid PK\\
  card\_id : string\\
  card\_type : string\\
  content : text\\
  position : integer\\
  message\_id : string\\
  conversation\_id : uuid FK\\
  inserted\_at, updated\_at
};

% predicted_codes
\node[tblsm, below right=0.8cm and 1.6cm of conv] (pred) {
  \textbf{predicted\_codes}
  \nodepart{two}
  \underline{id} : uuid PK\\
  code\_id : string\\
  cie10\_code : string\\
  confidence\_score : float\\
  status : string\\
  validated\_by, rejected\_by\\
  conversation\_id : uuid FK\\
  inserted\_at, updated\_at
};

% code_suggestions
\node[tblsm, below=1.0cm of conv] (sugg) {
  \textbf{code\_suggestions}
  \nodepart{two}
  \underline{id} : uuid PK\\
  suggestion\_id : string\\
  selected\_text : text\\
  suggested\_code : string\\
  suggested\_by : string\\
  conversation\_id : uuid FK\\
  inserted\_at, updated\_at
};

% cie10_codes
\node[tblref, right=1.6cm of pred] (cie10) {
  \textbf{cie10\_codes}
  \nodepart{two}
  \underline{id} : uuid PK\\
  code : string UNIQUE\\
  description : text\\
  type : string\\
  metadata : jsonb\\
  inserted\_at, updated\_at
};

%--- RELACIONES FK ---

% users 1 → N conversations
\draw[fk] (conv.west) -- node[above, font=\tiny\sffamily]{N} ++(-0.4,0)
  node[above, font=\tiny\sffamily, at end] {} -- (users.east)
  node[above, font=\tiny\sffamily, pos=0.9]{1};

% conversations 1 → N messages
\draw[fk] (msg.south) -- node[right, font=\tiny\sffamily]{N} ++(0,-0.3)
  -| (conv.north) node[left, font=\tiny\sffamily, pos=0.9]{1};

% conversations 1 → N analysis_cards
\draw[fk] (cards.west) -- node[above, font=\tiny\sffamily]{N} ++(-0.3,0)
  -- (conv.east) node[above, font=\tiny\sffamily, pos=0.1]{1};

% conversations 1 → N predicted_codes
\draw[fk] (pred.north) -- node[right, font=\tiny\sffamily]{N} ++(0,0.3)
  -| (conv.south) node[left, font=\tiny\sffamily, pos=0.9]{1};

% conversations 1 → N code_suggestions
\draw[fk] (sugg.north) -- node[right, font=\tiny\sffamily]{N} ++(0,0.2)
  -- (conv.south) node[right, font=\tiny\sffamily, pos=0.05]{1};

% predicted_codes N → 1 cie10_codes (referencia por código)
\draw[fk] (pred.east) -- node[above, font=\tiny\sffamily]{N} ++(0.3,0)
  -- (cie10.west) node[above, font=\tiny\sffamily, pos=0.9]{1};

\end{tikzpicture}
 dentro de figure+\centering+\resizebox{\textwidth}{!}{...}
%
% Tablas: users, conversations, messages, analysis_cards,
%         predicted_codes, code_suggestions, cie10_codes

\begin{tikzpicture}[
  every node/.style={font=\scriptsize\sffamily},
  %--- Estilo de tabla ERD (cabecera + campos) ---
  tbl/.style={
    rectangle split, rectangle split parts=2,
    rectangle split part fill={aquaESI!70, white},
    draw=aquaESI!80, thick, rounded corners=3pt,
    align=left, text width=3.8cm
  },
  tblsm/.style={
    rectangle split, rectangle split parts=2,
    rectangle split part fill={turquesa!60, white},
    draw=turquesa!70, thick, rounded corners=3pt,
    align=left, text width=3.4cm
  },
  tblref/.style={
    rectangle split, rectangle split parts=2,
    rectangle split part fill={gray!50, white},
    draw=gray!60, thick, rounded corners=3pt,
    align=left, text width=3.4cm
  },
  %--- Flechas de FK (1:N) ---
  fk/.style={->, >=Stealth, thick, gray!60},
  node distance=1.0cm and 1.6cm
]

%--- NODOS ---

% users
\node[tbl] (users) {
  \textbf{users}
  \nodepart{two}
  \underline{id} : uuid PK\\
  first\_name, last\_name\\
  username, email\\
  password\_hash\\
  locale, confirmed\_at\\
  inserted\_at, updated\_at
};

% conversations
\node[tbl, right=2.0cm of users] (conv) {
  \textbf{conversations}
  \nodepart{two}
  \underline{id} : uuid PK\\
  conversation\_id : string\\
  user\_id : uuid FK\\
  started\_at : utc\_datetime\\
  status : string\\
  deleted\_at : utc\_datetime\\
  inserted\_at, updated\_at
};

% messages
\node[tblsm, above right=0.8cm and 1.6cm of conv] (msg) {
  \textbf{messages}
  \nodepart{two}
  \underline{id} : uuid PK\\
  message\_id : string\\
  content : text\\
  user\_id, message\_type\\
  timestamp : utc\_datetime\\
  conversation\_id : uuid FK\\
  inserted\_at, updated\_at
};

% analysis_cards
\node[tblsm, right=1.6cm of conv] (cards) {
  \textbf{analysis\_cards}
  \nodepart{two}
  \underline{id} : uuid PK\\
  card\_id : string\\
  card\_type : string\\
  content : text\\
  position : integer\\
  message\_id : string\\
  conversation\_id : uuid FK\\
  inserted\_at, updated\_at
};

% predicted_codes
\node[tblsm, below right=0.8cm and 1.6cm of conv] (pred) {
  \textbf{predicted\_codes}
  \nodepart{two}
  \underline{id} : uuid PK\\
  code\_id : string\\
  cie10\_code : string\\
  confidence\_score : float\\
  status : string\\
  validated\_by, rejected\_by\\
  conversation\_id : uuid FK\\
  inserted\_at, updated\_at
};

% code_suggestions
\node[tblsm, below=1.0cm of conv] (sugg) {
  \textbf{code\_suggestions}
  \nodepart{two}
  \underline{id} : uuid PK\\
  suggestion\_id : string\\
  selected\_text : text\\
  suggested\_code : string\\
  suggested\_by : string\\
  conversation\_id : uuid FK\\
  inserted\_at, updated\_at
};

% cie10_codes
\node[tblref, right=1.6cm of pred] (cie10) {
  \textbf{cie10\_codes}
  \nodepart{two}
  \underline{id} : uuid PK\\
  code : string UNIQUE\\
  description : text\\
  type : string\\
  metadata : jsonb\\
  inserted\_at, updated\_at
};

%--- RELACIONES FK ---

% users 1 → N conversations
\draw[fk] (conv.west) -- node[above, font=\tiny\sffamily]{N} ++(-0.4,0)
  node[above, font=\tiny\sffamily, at end] {} -- (users.east)
  node[above, font=\tiny\sffamily, pos=0.9]{1};

% conversations 1 → N messages
\draw[fk] (msg.south) -- node[right, font=\tiny\sffamily]{N} ++(0,-0.3)
  -| (conv.north) node[left, font=\tiny\sffamily, pos=0.9]{1};

% conversations 1 → N analysis_cards
\draw[fk] (cards.west) -- node[above, font=\tiny\sffamily]{N} ++(-0.3,0)
  -- (conv.east) node[above, font=\tiny\sffamily, pos=0.1]{1};

% conversations 1 → N predicted_codes
\draw[fk] (pred.north) -- node[right, font=\tiny\sffamily]{N} ++(0,0.3)
  -| (conv.south) node[left, font=\tiny\sffamily, pos=0.9]{1};

% conversations 1 → N code_suggestions
\draw[fk] (sugg.north) -- node[right, font=\tiny\sffamily]{N} ++(0,0.2)
  -- (conv.south) node[right, font=\tiny\sffamily, pos=0.05]{1};

% predicted_codes N → 1 cie10_codes (referencia por código)
\draw[fk] (pred.east) -- node[above, font=\tiny\sffamily]{N} ++(0.3,0)
  -- (cie10.west) node[above, font=\tiny\sffamily, pos=0.9]{1};

\end{tikzpicture}
}
\caption{Diagrama entidad-relación de la base de datos de aplicación \texttt{cie10\_app}.}
\label{fig:erd-bd}
\end{figure}

\section{Tablas de la base de datos de aplicación (\texttt{cie10\_app})}

\subsection{Tabla \texttt{users}}

Almacena los usuarios registrados en el sistema.

\begin{table}[h]
\centering
\caption{Columnas de la tabla \texttt{users}.}
\label{tab:users}
\begin{tabular}{llll}
\hline
\textbf{Columna} & \textbf{Tipo} & \textbf{Restricciones} & \textbf{Descripción} \\
\hline
\texttt{id}                  & \texttt{uuid}          & PK, NOT NULL      & Identificador único. \\
\texttt{first\_name}         & \texttt{varchar}       & NOT NULL          & Nombre del usuario. \\
\texttt{last\_name}          & \texttt{varchar}       & NOT NULL          & Apellidos. \\
\texttt{username}            & \texttt{varchar}       & UNIQUE, NOT NULL  & Nombre de usuario. \\
\texttt{email}               & \texttt{varchar}       & UNIQUE, NOT NULL  & Correo electrónico. \\
\texttt{password\_hash}      & \texttt{varchar}       & NOT NULL          & Hash bcrypt de la contraseña. \\
\texttt{confirmation\_token} & \texttt{varchar}       & NULLABLE          & Token de confirmación por email. \\
\texttt{confirmed\_at}       & \texttt{timestamptz}   & NULLABLE          & Fecha de confirmación de la cuenta. \\
\texttt{locale}              & \texttt{varchar}       & DEFAULT \texttt{'es'} & Idioma preferido (\texttt{es} o \texttt{en}). \\
\texttt{inserted\_at}        & \texttt{timestamptz}   & NOT NULL          & Fecha de creación. \\
\texttt{updated\_at}         & \texttt{timestamptz}   & NOT NULL          & Fecha de última modificación. \\
\hline
\end{tabular}
\end{table}

\subsection{Tabla \texttt{conversations}}

Proyección CQRS de las conversaciones (sesiones de análisis) de cada usuario.

\begin{table}[h]
\centering
\caption{Columnas de la tabla \texttt{conversations}.}
\label{tab:conversations}
\begin{tabular}{llll}
\hline
\textbf{Columna} & \textbf{Tipo} & \textbf{Restricciones} & \textbf{Descripción} \\
\hline
\texttt{id}               & \texttt{uuid}        & PK, NOT NULL        & Identificador interno. \\
\texttt{conversation\_id} & \texttt{varchar}     & UNIQUE, NOT NULL    & UUID generado por el cliente. \\
\texttt{user\_id}         & \texttt{varchar}     & NOT NULL, INDEX     & Referencia al usuario propietario. \\
\texttt{started\_at}      & \texttt{timestamptz} & NOT NULL            & Momento de inicio de la conversación. \\
\texttt{status}           & \texttt{varchar}     & DEFAULT \texttt{'active'} & Estado: \texttt{active}. \\
\texttt{deleted\_at}      & \texttt{timestamptz} & NULLABLE, INDEX     & Borrado lógico (soft delete). \\
\texttt{inserted\_at}     & \texttt{timestamptz} & NOT NULL            & Fecha de inserción en la proyección. \\
\texttt{updated\_at}      & \texttt{timestamptz} & NOT NULL            & Fecha de última actualización. \\
\hline
\end{tabular}
\end{table}

\textbf{Índices}: \texttt{(user\_id)}, \texttt{(deleted\_at)}.

\subsection{Tabla \texttt{messages}}

Proyección CQRS de los mensajes intercambiados dentro de cada conversación.

\begin{table}[h]
\centering
\caption{Columnas de la tabla \texttt{messages}.}
\label{tab:messages}
\begin{tabular}{llll}
\hline
\textbf{Columna} & \textbf{Tipo} & \textbf{Restricciones} & \textbf{Descripción} \\
\hline
\texttt{id}               & \texttt{uuid}        & PK, NOT NULL   & Identificador interno. \\
\texttt{message\_id}      & \texttt{varchar}     & UNIQUE, NOT NULL & UUID del mensaje. \\
\texttt{content}          & \texttt{text}        & NOT NULL       & Contenido textual del mensaje. \\
\texttt{user\_id}         & \texttt{varchar}     & NOT NULL       & Usuario que emitió el mensaje. \\
\texttt{timestamp}        & \texttt{timestamptz} & NOT NULL, INDEX & Momento de emisión. \\
\texttt{message\_type}    & \texttt{varchar}     & NOT NULL       & Tipo: \texttt{user}, \texttt{system} o \texttt{analysis}. \\
\texttt{conversation\_id} & \texttt{uuid}        & FK → \texttt{conversations}, CASCADE & Conversación a la que pertenece. \\
\texttt{inserted\_at}     & \texttt{timestamptz} & NOT NULL       & Fecha de inserción. \\
\texttt{updated\_at}      & \texttt{timestamptz} & NOT NULL       & Fecha de modificación. \\
\hline
\end{tabular}
\end{table}

\textbf{Índices}: \texttt{(conversation\_id)}, \texttt{(message\_id)}, \texttt{(timestamp)}.

\subsection{Tabla \texttt{analysis\_cards}}

Proyección CQRS de las tarjetas de respuesta generadas por el motor de IA para cada mensaje de análisis.

\begin{table}[h]
\centering
\caption{Columnas de la tabla \texttt{analysis\_cards}.}
\label{tab:cards}
\begin{tabular}{llll}
\hline
\textbf{Columna} & \textbf{Tipo} & \textbf{Restricciones} & \textbf{Descripción} \\
\hline
\texttt{id}               & \texttt{uuid}    & PK, NOT NULL           & Identificador interno. \\
\texttt{card\_id}         & \texttt{varchar} & UNIQUE, NOT NULL       & UUID de la tarjeta. \\
\texttt{card\_type}       & \texttt{varchar} & NOT NULL               & Tipo: \texttt{summary}, \texttt{codes} o \texttt{recommendations}. \\
\texttt{content}          & \texttt{text}    & NOT NULL               & JSON (tipo \texttt{codes}) o texto libre. \\
\texttt{position}         & \texttt{integer} & NOT NULL               & Orden de renderizado dentro del mensaje. \\
\texttt{message\_id}      & \texttt{varchar} & NOT NULL, INDEX        & UUID del mensaje al que pertenece. \\
\texttt{conversation\_id} & \texttt{uuid}    & FK → \texttt{conversations}, CASCADE & Conversación. \\
\texttt{inserted\_at}     & \texttt{timestamptz} & NOT NULL           & Fecha de inserción. \\
\texttt{updated\_at}      & \texttt{timestamptz} & NOT NULL           & Fecha de modificación. \\
\hline
\end{tabular}
\end{table}

\textbf{Índices}: \texttt{(conversation\_id)}, \texttt{(message\_id)}.

\subsection{Tabla \texttt{predicted\_codes}}

Proyección CQRS de los códigos CIE-10-ES predichos por el motor de IA, con su ciclo de vida de validación.

\begin{table}[h]
\centering
\caption{Columnas de la tabla \texttt{predicted\_codes}.}
\label{tab:predicted-codes}
\begin{tabular}{llll}
\hline
\textbf{Columna} & \textbf{Tipo} & \textbf{Restricciones} & \textbf{Descripción} \\
\hline
\texttt{id}               & \texttt{uuid}    & PK, NOT NULL           & Identificador interno. \\
\texttt{code\_id}         & \texttt{varchar} & UNIQUE, NOT NULL       & UUID de la predicción. \\
\texttt{cie10\_code}      & \texttt{varchar} & NOT NULL, INDEX        & Código CIE-10-ES (p. ej., \texttt{E11}). \\
\texttt{reasoning}        & \texttt{text}    & NOT NULL               & Explicación textual del modelo. \\
\texttt{confidence\_score}& \texttt{float}   & NOT NULL               & Puntuación de confianza [0, 1]. \\
\texttt{status}           & \texttt{varchar} & DEFAULT \texttt{'pending'}, INDEX & Estado: \texttt{pending}, \texttt{validated} o \texttt{rejected}. \\
\texttt{validated\_by}    & \texttt{varchar} & NULLABLE               & ID del usuario que validó. \\
\texttt{rejected\_by}     & \texttt{varchar} & NULLABLE               & ID del usuario que rechazó. \\
\texttt{rejection\_reason}& \texttt{text}    & NULLABLE               & Motivo del rechazo. \\
\texttt{conversation\_id} & \texttt{uuid}    & FK → \texttt{conversations}, CASCADE & Conversación. \\
\texttt{inserted\_at}     & \texttt{timestamptz} & NOT NULL           & Fecha de inserción. \\
\texttt{updated\_at}      & \texttt{timestamptz} & NOT NULL           & Fecha de modificación. \\
\hline
\end{tabular}
\end{table}

\textbf{Índices}: \texttt{(conversation\_id)}, \texttt{(code\_id)}, \texttt{(cie10\_code)}, \texttt{(status)}.

\subsection{Tabla \texttt{cie10\_codes}}

Catálogo de referencia de los códigos CIE-10-ES cargados desde los ficheros CSV obtenidos del portal eCIE-maps.

\begin{table}[h]
\centering
\caption{Columnas de la tabla \texttt{cie10\_codes}.}
\label{tab:cie10codes}
\begin{tabular}{llll}
\hline
\textbf{Columna} & \textbf{Tipo} & \textbf{Restricciones} & \textbf{Descripción} \\
\hline
\texttt{id}          & \texttt{uuid}    & PK, NOT NULL       & Identificador interno. \\
\texttt{code}        & \texttt{varchar} & UNIQUE, NOT NULL   & Código alfanumérico (p. ej., \texttt{E11}, \texttt{I10}). \\
\texttt{description} & \texttt{text}    & NOT NULL           & Descripción oficial en español. \\
\texttt{type}        & \texttt{varchar} & NOT NULL, INDEX    & Categoría: \texttt{diagnosis}, \texttt{procedure} o \texttt{chemical}. \\
\texttt{metadata}    & \texttt{jsonb}   & NULLABLE           & Metadatos clínicos (capítulo, restricciones demográficas, etc.). \\
\texttt{inserted\_at}& \texttt{timestamptz} & NOT NULL       & Fecha de carga. \\
\texttt{updated\_at} & \texttt{timestamptz} & NOT NULL       & Fecha de actualización. \\
\hline
\end{tabular}
\end{table}

El catálogo contiene aproximadamente 101.246 diagnósticos, 78.496 procedimientos y 5.050 sustancias químicas.

\subsection{Tabla \texttt{code\_suggestions}}

Sugerencias manuales de código realizadas por el usuario al seleccionar un fragmento de texto del informe.

\begin{table}[h]
\centering
\caption{Columnas de la tabla \texttt{code\_suggestions}.}
\label{tab:suggestions}
\begin{tabular}{llll}
\hline
\textbf{Columna} & \textbf{Tipo} & \textbf{Restricciones} & \textbf{Descripción} \\
\hline
\texttt{id}               & \texttt{uuid}    & PK, NOT NULL     & Identificador interno. \\
\texttt{suggestion\_id}   & \texttt{varchar} & UNIQUE, NOT NULL & UUID de la sugerencia. \\
\texttt{selected\_text}   & \texttt{varchar} & NOT NULL         & Fragmento de texto seleccionado. \\
\texttt{suggested\_code}  & \texttt{varchar} & NOT NULL         & Código CIE-10-ES propuesto (mayúsculas). \\
\texttt{suggested\_by}    & \texttt{varchar} & NOT NULL         & ID del usuario que realizó la sugerencia. \\
\texttt{conversation\_id} & \texttt{uuid}    & FK → \texttt{conversations}, CASCADE & Conversación. \\
\texttt{inserted\_at}     & \texttt{timestamptz} & NOT NULL     & Fecha de creación. \\
\texttt{updated\_at}      & \texttt{timestamptz} & NOT NULL     & Fecha de modificación. \\
\hline
\end{tabular}
\end{table}

\subsection{Tabla \texttt{projection\_versions}}

Tabla de control interno de Commanded para garantizar la coherencia de las proyecciones al reproducir el log de eventos.

\begin{table}[h]
\centering
\caption{Columnas de la tabla \texttt{projection\_versions}.}
\label{tab:proj-versions}
\begin{tabular}{llll}
\hline
\textbf{Columna} & \textbf{Tipo} & \textbf{Restricciones} & \textbf{Descripción} \\
\hline
\texttt{projection\_name}         & \texttt{varchar} & PK          & Nombre del proyector. \\
\texttt{last\_seen\_event\_number}& \texttt{bigint}  & NOT NULL    & Último evento procesado por este proyector. \\
\texttt{inserted\_at}             & \texttt{timestamptz} & NOT NULL & Fecha de creación. \\
\texttt{updated\_at}              & \texttt{timestamptz} & NOT NULL & Fecha de actualización. \\
\hline
\end{tabular}
\end{table}

\section{Base de datos de eventos (\texttt{cie10\_eventstore})}

La base de datos \texttt{cie10\_eventstore} es gestionada íntegramente por la biblioteca \texttt{EventStore} (dependencia de Commanded). Su esquema incluye las tablas \texttt{streams}, \texttt{events} y \texttt{subscriptions}, que almacenan el log de eventos inmutable. No se modifica directamente; toda la interacción se realiza a través del módulo \texttt{App.CommandedApplication}.

Los eventos registrados son: \texttt{ConversationStarted}, \texttt{MessageSent}, \texttt{AnalysisRequested}, \texttt{AIPredictionReceived}, \texttt{CodeValidated} y \texttt{CodeRejected}. Cada evento incluye los campos de la acción más metadatos de Commanded (\texttt{causation\_id}, \texttt{correlation\_id}, \texttt{created\_at}).

\section{Relaciones entre entidades}

\begin{itemize}
    \item Un \textbf{usuario} tiene cero o más \textbf{conversaciones}.
    \item Una \textbf{conversación} tiene cero o más \textbf{mensajes}, \textbf{tarjetas de análisis}, \textbf{códigos predichos} y \textbf{sugerencias de código}. Al eliminar una conversación, todos sus registros dependientes se eliminan en cascada.
    \item Un \textbf{mensaje} de tipo \texttt{analysis} tiene exactamente tres \textbf{tarjetas de análisis}: \texttt{summary} (posición 0), \texttt{codes} (posición 1) y \texttt{recommendations} (posición 2).
    \item Un \textbf{código predicho} hace referencia a un código del \textbf{catálogo CIE-10-ES} a través del campo \texttt{cie10\_code}, aunque no existe una clave foránea explícita para permitir predicciones de códigos aún no importados al catálogo.
\end{itemize}
